% Figure: map of the analysis (Irene: "use graphs to help readers understand the
% analysis"). Drawn in TikZ so it uses the report's fonts. Original diagram.
% Introduced by Pablo at the start of Section 3 and referred to in the
% Introduction road map. Numbers checked against results/ on 22 Nov 2026.
\begin{figure*}[t]
\centering
\begin{tikzpicture}[
  font=\footnotesize,
  fdata/.style={draw=FIELDdata, fill=FIELDdata!8, rounded corners=2pt, align=center, text width=3.75cm, minimum height=0.9cm, inner sep=3pt},
  fstep/.style={draw=FIELDstep, fill=FIELDstep!7, rounded corners=2pt, align=left, text width=7.55cm, minimum height=0.55cm, inner sep=3pt},
  fgoal/.style={draw=black!45, fill=black!3, rounded corners=2pt, align=left, text width=4.6cm, minimum height=0.55cm, inner sep=3pt},
  arr/.style={-{Stealth[length=4pt]}, draw=black!55, line width=0.5pt},
  hdr/.style={font=\footnotesize\bfseries, anchor=south}
]
% columns
\node[hdr] at (0,0.35) {Data (\S\ref{sec:data})};
\node[hdr] at (6.55,0.35) {Analysis (\S\ref{sec:methodology})};
\node[hdr] at (12.95,0.35) {Question answered (\S\ref{sec:results})};
% data
\node[fdata] (es) at (0,-1.6) {\textbf{Spain} (Ministry, SIIU)\\12,125 field means\\89 universities, 2015--23};
\node[fdata] (pt) at (0,-4.0) {\textbf{Portugal} (UCI)\\3,703 students, 15 degrees};
\node[fdata] (us) at (0,-5.25) {\textbf{United States}\\33 universities, 1,845 cells};
% steps
\node[fstep] (s1) at (6.55,0)    {Deviations, decomposition, mixed model (eqs.~\ref{eq:deviation}--\ref{eq:mixed})};
\node[fstep] (s2) at (6.55,-0.8) {Five models: tuned 2018--21, tested 2022--23 (Alg.~\ref{alg:tuning})};
\node[fstep] (s3) at (6.55,-1.6) {Controls for intake, size and dropout};
\node[fstep] (s4) at (6.55,-2.4) {Top-20\% selection simulation (eq.~\ref{eq:selection})};
\node[fstep] (s5) at (6.55,-3.2) {Adjusted marks from past deviations (eq.~\ref{eq:adjusted})};
\node[fstep] (s6) at (6.55,-4.2) {Four ranking rules judged on the next semester};
\node[fstep] (s7) at (6.55,-5.25) {Same model on ten broad areas in both countries};
% goals
\node[fgoal] (g1) at (12.95,0)    {O1 How much is field? (\S\ref{sec:res-field})};
\node[fgoal] (g2) at (12.95,-0.8) {O2 Persists and predicts? (\S\ref{sec:res-prediction})};
\node[fgoal] (g3) at (12.95,-1.6) {O3 Explained by intake? (\S\ref{sec:res-confounders})};
\node[fgoal] (g4) at (12.95,-2.4) {O4 Who gets selected? (\S\ref{sec:res-selection})};
\node[fgoal] (g5) at (12.95,-3.2) {O5 Does a correction fix it? (\S\ref{sec:res-adjusted})};
\node[fgoal] (g6) at (12.95,-4.2) {O5 At what cost? (\S\ref{sec:res-rules})};
\node[fgoal] (g7) at (12.95,-5.25) {O6 Same in the US? (\S\ref{sec:res-us})};
% arrows data -> steps
\foreach \s in {s1,s2,s3,s4,s5} \draw[arr] (es.east) -- (\s.west);
\draw[arr] (pt.east) -- (s3.west);
\draw[arr] (pt.east) -- (s6.west);
\draw[arr] (us.east) -- (s7.west);
\draw[arr] (es.east) to[out=-30,in=180] (s7.west);
% arrows steps -> goals
\foreach \i in {1,...,7} \draw[arr] (s\i.east) -- (g\i.west);
\end{tikzpicture}
\caption{Map of our analysis: datasets on the left feed the methods in the centre, which answer the numbered questions and results subsections on the right. We analysed the datasets separately without linking individual records between countries. Arrows show analytical use, not data linkage; displayed sample counts come from the saved results. Original diagram.}
\label{fig:overview}
\end{figure*}
